\documentclass[11pt,letterpaper]{article}
\usepackage[margin=1in]{geometry}
\usepackage{booktabs}
\usepackage{tabularx}
\usepackage{longtable}
\usepackage{array}
\usepackage{hyperref}
\usepackage{amsmath}
\usepackage{parskip}
\usepackage{titlesec}
\usepackage{enumitem}
\usepackage[numbers]{natbib}

\titlespacing*{\section}{0pt}{1.2ex plus .2ex}{0.8ex}
\titlespacing*{\subsection}{0pt}{1ex plus .2ex}{0.5ex}

\hypersetup{colorlinks=true, linkcolor=blue, citecolor=blue, urlcolor=blue}

\title{\textbf{AI Virtual Cell Models for Dopaminergic Neuron Perturbation Prediction}\\
\large Midterm Progress Report}
\author{Katherine Deborah Godwin Gnanaraj \and Zihan Jin \and Yumejichi Fujita}
\date{March 2026 \quad CSCI 566}

\begin{document}
\maketitle

%% ─────────────────────────────────────────────────────────────────────────────
\begin{abstract}
Parkinson's disease (PD) is driven by the selective loss of dopaminergic (DA) neurons in the
substantia nigra. Identifying which gene perturbations can shift cells toward or away from a DA
neuron transcriptional state is critical for discovering therapeutic targets, but testing hundreds
of candidates experimentally is prohibitively costly. AI virtual cell (AIVC) models offer a
computational alternative: trained on single-cell CRISPR perturbation screens, these models
learn to predict post-perturbation gene expression profiles. In this midterm report, we survey
the landscape of AIVC perturbation prediction models, describe our datasets and preprocessing
pipeline, and present benchmarking results on the GSE152988 iPSC-derived neuron dataset. We
evaluate a mean expression baseline, scGen, and GEARS (Katherine Deborah Godwin Gnanaraj), a
20-epoch full-gene GEARS variant (Yumejichi Fujita), and STATE/CPA fine-tuned on a combined
CRISPRi+CRISPRa dataset with a zero-shot cross-dataset evaluation design (Zihan Jin). All
models are evaluated under both standard absolute metrics and the stricter delta-based
perturbation-specific metrics we designed as our primary evaluation framework. GEARS achieves a
top-20 DEG delta Pearson of 0.262 (CRISPRi) on completely unseen gene knockdowns; Yumejichi
Fujita's 20-epoch run achieves Pearson DE of 0.968 on an independent split using all 33k genes.
We outline a roadmap for integrating STATE/CPA results and extending to DA identity scoring.
\end{abstract}

%% ─────────────────────────────────────────────────────────────────────────────
\section{Introduction}

Parkinson's disease affects approximately 10 million people worldwide and currently has no
disease-modifying treatment. The primary pathological event is the degeneration of dopaminergic
(DA) neurons in the substantia nigra pars compacta (SNpc), which produce dopamine essential for
motor control. Understanding the transcriptional programs that define DA neuron identity --- and
how genetic perturbations shift cells toward or away from this state --- is a central goal of
modern PD research.

CRISPR-based perturbation screens, combined with single-cell RNA sequencing (scRNA-seq), now
make it possible to measure how hundreds of gene knockdowns reshape transcriptional profiles at
single-cell resolution. However, systematically testing all $\sim$20,000 human protein-coding
genes experimentally remains infeasible. AI virtual cell (AIVC) models address this by learning
to predict post-perturbation gene expression from observed data, enabling in silico screening of
genetic targets at scale.

In this project, we benchmark several AIVC perturbation prediction models on two human datasets
featuring iPSC-derived neurons subjected to CRISPR perturbation screens. Our primary goal is to
identify which gene perturbations most strongly promote or inhibit a DA neuron transcriptional
state, contributing to a systematic gene ranking pipeline for experimental validation.

This midterm report covers: (1) a literature review of AIVC models, (2) dataset description and
preprocessing, (3) initial experimental results from the mean baseline, scGen, and GEARS on
GSE152988, evaluated under our full metric suite, and (4) a research roadmap for the remainder
of the project.

%% ─────────────────────────────────────────────────────────────────────────────
\section{Related Work}

A growing ecosystem of AI models has been developed for single-cell perturbation prediction.
We survey the most relevant approaches below and summarize their key characteristics in
Table~\ref{tab:models}.

\textbf{scGen}~\cite{lopez2018} is a variational autoencoder (VAE) that models perturbation
effects via latent space arithmetic. Control and perturbed cells are encoded into a shared latent
space; a perturbation ``delta vector''---the difference between mean latent representations of
perturbed and control cells---is applied to new control cells to predict their perturbed state.
scGen is conceptually simple and fast to train, but cannot generalize to novel perturbations not
seen during training.

\textbf{GEARS}~\cite{roohani2023} overcomes this limitation by incorporating a gene co-expression
graph (derived from the STRING database) as an inductive bias. A graph neural network (GNN) is
trained where each node is a gene and edges encode co-expression relationships. By reasoning over
this gene network, GEARS can predict the effect of genes never perturbed in training. It is
currently the standard baseline for unseen perturbation evaluation.

\textbf{CPA}~\cite{lotfollahi2021} extends the VAE framework by learning disentangled latent
representations for cell identity, perturbation effect, and technical covariates, enabling
compositional prediction of perturbation combinations.

\textbf{STATE} is a perturbation prediction framework available through the CZ AIVC portal. It
supports model backends including STATE, CPA, scVI, scGPT, and Tahoe. We fine-tune the
\texttt{ST-SE-Replogle} pretrained model on GSE152988 via this framework (Zihan Jin).

\textbf{scFoundation}~\cite{hao2024} is a large-scale foundation model pretrained on over 50
million human single-cell transcriptomes using a masked gene expression modeling objective over a
fixed 19,264-gene vocabulary. It acts as a feature extractor generating context-aware gene
embeddings (cells $\times$ 19,264 $\times$ 512) for downstream models such as GEARS. Requires
40+ GB VRAM; planned for HPC or Colab A100.

\textbf{scGenePT} (2024) augments scGen with gene-level text embeddings from biological
literature, improving generalization to unseen perturbations.

\textbf{scGPT}~\cite{cui2024} is a 100M-parameter transformer pretrained on 33 million human
single-cell profiles, supporting perturbation prediction via fine-tuning. Requires HPC.

\begin{table}[h]
\centering
\caption{Summary of AIVC perturbation prediction models surveyed and evaluated.}
\label{tab:models}
\small
\begin{tabular}{lllllll}
\toprule
Model & Architecture & Unseen & Gene Graph & Foundation & Owner & Status \\
\midrule
Mean Baseline            & ---                   & No      & No      & No              & KG  & Done \\
scGen                    & VAE                   & No      & No      & No              & KG  & Done \\
GEARS (5 ep, 5k HVG)     & GNN                   & Yes     & STRING  & No              & KG  & Done \\
GEARS (20 ep, all genes) & GNN                   & Yes     & None$^*$& No              & YF  & Done \\
GEARS (20 ep, Colab)     & GNN                   & Yes     & STRING  & No              & KG  & Done \\
STATE (GSE152988 comb.)  & Struct.+CPA/scVI      & Partial & No      & Replogle        & ZJ  & In progress \\
STATE (GSE124703)        & Structured            & Partial & No      & Replogle        & ZJ  & In progress \\
scFoundation+GEARS       & GNN+Transformer       & Yes     & STRING  & Yes (50M cells) & --- & Planned \\
scGenePT                 & VAE+LLM text          & Partial & No      & Text embeddings & --- & Planned \\
scGPT                    & Transformer           & Yes     & No      & Yes (33M cells) & --- & Planned \\
\bottomrule
\end{tabular}
\vspace{2pt}
{\footnotesize $^*$Empty co-expression graph used due to library compatibility issue with full gene set.}
\end{table}

%% ─────────────────────────────────────────────────────────────────────────────
\section{Datasets}

\subsection{GSE152988 --- Primary Benchmark Dataset}

Our primary benchmark is the Tian \& Kampmann (2021) iPSC-derived neuron CRISPR screen~\cite{tian2021},
accessed as pre-processed AnnData objects. The dataset contains two perturbation modalities:

\begin{itemize}[noitemsep]
  \item \textbf{CRISPRi (transcriptional repression):} 32,300 cells across 184 gene knockdowns
        + 437 control cells; 33,538 genes measured.
  \item \textbf{CRISPRa (transcriptional activation):} 21,193 cells across 100 gene activations
        + 434 control cells; 33,538 genes measured.
\end{itemize}

Both datasets feature iPSC-induced neurons from the LUHMES cell line. Raw integer count matrices
were normalized to 10,000 counts per cell and log1p-transformed. For GEARS experiments on local
hardware (NVIDIA GTX 1650 Ti, 4 GB VRAM), the top 5,000 highly variable genes (HVGs) were
selected using the Seurat v3 method to reduce memory requirements. For Colab (A100) experiments,
all 33,538 genes are retained: Yumejichi Fujita confirmed that full-gene evaluation is tractable
on Colab hardware and produces similar Pearson DE results (0.968 vs.\ 0.973 at 5k HVGs). DA
marker gene analysis (71 canonical markers including TH, SLC6A3, NR4A2, FOXA2, PITX3) reveals
that 67/71 markers are present in the gene panel.

\subsection{GSE124703 --- Secondary Dataset / Zero-Shot Transfer Target}

This dataset~\cite{tian2019} contains iPSC and Day 7 neuron cells with 57 CRISPRi guide
sequences targeting $\sim$25 gene knockdowns, providing two developmental contexts (iPSC vs.\
early neuron stage) for the same perturbations. Per our two-track evaluation design
(Section~\ref{sec:methods}), GSE124703 is held out from training until the zero-shot benchmark
is completed.

\subsection{GSE140231 --- DA Reference Dataset}

This substantia nigra dataset ($\sim$20,000 cells from human donors) contains confirmed DA
neurons alongside other cell types. We plan to use it to derive clean DA neuron identity labels
for scoring perturbation effects in terms of DA identity shift. Preprocessing is ongoing.

%% ─────────────────────────────────────────────────────────────────────────────
\section{Methods}
\label{sec:methods}

\subsection{Preprocessing}

Raw count matrices were normalized to 10,000 counts per cell and log1p-transformed using
Scanpy. For GEARS, gene names were mapped to GEARS' format (\texttt{GENE+ctrl} for single
perturbations, \texttt{ctrl} for controls). For the combined dataset (STATE/CPA, Zihan Jin),
CRISPRi and CRISPRa cells are merged with modality-specific labels
(\texttt{CRISPRI::GENE}, \texttt{CRISPRA::GENE}) and modality-prefixed batch labels.

\subsection{Two-Track Evaluation Design}

\textbf{Track A --- Clean Zero-Shot Cross-Dataset Evaluation (STATE):}
Fine-tune \texttt{ST-SE-Replogle} on GSE152988 (CRISPRi + CRISPRa combined). Evaluate zero-shot
on GSE124703 (completely unseen dataset, unseen cell developmental stages). GSE124703 must not
be used in training until after this benchmark.

\textbf{Track B --- Perturbation Generalization Within Dataset (GEARS, scGen):}
Hold out 10\% of perturbation conditions from GSE152988 entirely from training. The model must
predict effects of genes never perturbed in training, relying on the gene co-expression graph
(GEARS) or seen-perturbation interpolation (scGen).

\subsection{Data Splits}

For Track B (GEARS, scGen): \texttt{zero\_shot\_pert} split --- perturbations with $\geq$30
cells split 80/10/10 (train/val/test) by condition, seed = 42. CRISPRi: 148/18/18; CRISPRa:
80/10/10. scGen is additionally evaluated on 20 randomly selected seen perturbations (80\% cells
train, 20\% held out per condition) --- a strictly easier task, not directly comparable to
GEARS' unseen numbers. For Track A (STATE): 228/28/28 perturbations on the combined dataset.

\subsection{Evaluation Metric Suite}

Following the team evaluation framework, we use a two-tier metric suite:

\textbf{Primary --- delta-based (perturbation-specific):} For each held-out perturbation, compute
$\overline{x}_{\text{ctrl}}$, $\text{true\_delta} = \overline{x}_{\text{pert}} - \overline{x}_{\text{ctrl}}$,
$\text{pred\_delta} = \hat{x}_{\text{pert}} - \overline{x}_{\text{ctrl}}$, then report:
\begin{itemize}[noitemsep]
  \item \textbf{Pearson $\Delta$ (top-20 DEGs):} Correlation of predicted vs.\ true change on
        the 20 genes most differentially expressed vs.\ control
  \item \textbf{Pearson $\Delta$ (DA markers):} Same, restricted to the 71-gene DA marker panel
  \item \textbf{Sign accuracy (top-20 DEGs / DA markers):} Fraction with correct up/down direction
\end{itemize}

\textbf{Secondary --- absolute:} Pearson $r$ (all genes), Pearson $r$ (top-20 DEGs), Pearson $r$
(DA markers), MSE.

\textbf{Why this matters:} Single-gene perturbations affect only a small fraction of the
transcriptome. Cell identity explains far more variance than any perturbation, so absolute Pearson
across all genes will exceed 0.99 even for a trivial mean baseline~\cite{boiarsky2023,wenk2024}.
The delta-based metrics isolate the perturbation-specific signal.

\subsection{Data Leakage Considerations}

For any pretrained model used in this project (scFoundation pretrained on $\sim$50M cells from
CellxGene, STATE pretrained on \texttt{ST-SE-Replogle}), GSE152988 or GSE124703 may already be
present in the pretraining corpus. We flag this as an open question requiring verification before
drawing strong conclusions about the benefit of foundation model initialization.

%% ─────────────────────────────────────────────────────────────────────────────
\section{Results}

\subsection{Evaluation Protocol Differences}

The three model families use different evaluation protocols and \textbf{cannot be ranked directly}:

\begin{table}[h]
\centering
\small
\begin{tabular}{lllp{5.5cm}}
\toprule
Model & Split Type & \# Test Perts & Protocol \\
\midrule
Mean Baseline      & Seen            & 20    & 20 conditions; predict with global mean \\
scGen              & Seen            & 20    & 20\% of cells per condition held out \\
GEARS (5 ep)       & \textbf{Unseen} & 46    & Entire conditions held out (simulation split) \\
GEARS (20 ep, Colab) & \textbf{Unseen} & $\sim$18 & Entire conditions held out (zero\_shot\_pert) \\
\bottomrule
\end{tabular}
\end{table}

\subsection{CRISPRi Results}

\begin{table}[h]
\centering
\caption{CRISPRi --- primary delta-based metrics.}
\small
\begin{tabular}{lcccccc}
\toprule
Model & Pearson $\Delta$ top-20 & Pearson $\Delta$ DA & Sign top-20 & Sign DA & \# Perts & Protocol \\
\midrule
Mean Baseline        & 0.769            & 0.518 & 0.953 & 0.653 & 20 & Seen (cond.\ holdout) \\
scGen                & 0.442            & 0.168 & 0.738 & 0.478 & 20 & Seen (80/20 cell) \\
GEARS (5 ep, 5k HVG) & 0.262$^\dagger$ & ---   & 0.663$^\dagger$ & --- & 46 & \textbf{Unseen} \\
GEARS (20 ep, Colab) & 0.129$^\dagger$ & ---   & 0.623$^\dagger$ & --- & 37 & \textbf{Unseen} \\
\bottomrule
\end{tabular}
\vspace{2pt}
\begin{minipage}{\linewidth}
{\footnotesize $^\dagger$From GEARS' own internal test evaluation (computed against GEARS' stored ctrl mean).
DA marker delta columns omitted: our custom evaluation's ctrl\_mean is incompatible with GEARS' internal
normalization space for CRISPRi, yielding near-zero deltas. Absolute Pearson on DA markers (Table~\ref{tab:crispri_abs})
is unaffected and reliable. 58/66 DA markers in panel.}
\end{minipage}
\end{table}

\begin{table}[h]
\centering
\caption{CRISPRi --- secondary absolute metrics (for cross-paper comparison).}
\label{tab:crispri_abs}
\small
\begin{tabular}{lccccc}
\toprule
Model & Pearson $r$ (all) & Pearson $r$ (top-20) & Pearson $r$ (DA) & MSE & Protocol \\
\midrule
Mean Baseline        & 0.9988 & 0.9917          & 0.9998 & 0.0005  & Seen (20) \\
scGen                & 0.9938 & 0.9523          & 0.9988 & 0.0025  & Seen (20) \\
GEARS (5 ep)         & 0.9965 & 0.9733          & ---    & 0.00137 & \textbf{Unseen (46)} \\
GEARS (20 ep, Colab) & 0.9963 & \textbf{0.9776} & 0.9984 & 0.0007 & \textbf{Unseen (37)} \\
\bottomrule
\end{tabular}
\end{table}

All methods achieve Pearson $r > 0.99$ across all genes, confirming this metric is dominated by
cell-type identity. The mean baseline's high delta Pearson (0.769) is an artifact of evaluation
design: randomly chosen seen perturbations are representative of the dataset average, consistent
with the literature~\cite{boiarsky2023,wenk2024}. scGen underperforms the mean baseline on delta
metrics (0.442 vs.\ 0.769) despite having 80\% of eval perturbation cells in training, suggesting
its latent space arithmetic introduces noise in the delta direction. GEARS' \textbf{sign accuracy
of 62.3\%} on completely unseen perturbations is above the 50\% chance baseline. The 20-epoch
Colab model achieves lower Pearson delta (0.129) than the 5-epoch run (0.262) due to the harder
test set (37 perts, stricter simulation selection), not model degradation --- it achieves higher
Pearson DE (0.978 vs.\ 0.973).

\subsection{CRISPRa Results}

\begin{table}[h]
\centering
\caption{CRISPRa --- primary delta-based metrics.}
\small
\begin{tabular}{lcccccc}
\toprule
Model & Pearson $\Delta$ top-20 & Pearson $\Delta$ DA & Sign top-20 & Sign DA & \# Perts & Protocol \\
\midrule
Mean Baseline        & 0.591           & 0.628           & 0.920 & 0.493 & 20 & Seen (cond.\ holdout) \\
scGen                & 0.591           & 0.283           & 0.833 & 0.407 & 20 & Seen (80/20 cell) \\
GEARS (5 ep)         & 0.334$^\dagger$ & ---             & 0.714$^\dagger$ & --- & 25 & \textbf{Unseen} \\
GEARS (20 ep, Colab) & \textbf{0.398}  & \textbf{0.498}  & \textbf{0.793} & 0.510 & 20 & \textbf{Unseen} \\
\bottomrule
\end{tabular}
\vspace{2pt}
{\footnotesize $^\dagger$From GEARS' own internal test evaluation. GEARS (20 ep) delta metrics from custom evaluation loop with ctrl\_mean from GEARS train\_loader. 58/66 DA markers in 33k gene panel.}
\end{table}

\begin{table}[h]
\centering
\caption{CRISPRa --- secondary absolute metrics.}
\small
\begin{tabular}{lccccc}
\toprule
Model & Pearson $r$ (all) & Pearson $r$ (top-20) & Pearson $r$ (DA) & MSE & Protocol \\
\midrule
Mean Baseline        & 0.9956 & 0.9738          & 0.9964 & 0.0018  & Seen (20) \\
scGen                & 0.9928 & 0.9582          & 0.9960 & 0.0034  & Seen (20) \\
GEARS (5 ep)         & 0.9948 & 0.9565          & ---    & 0.00224 & \textbf{Unseen (25)} \\
GEARS (20 ep, Colab) & 0.9941 & \textbf{0.9614} & 0.9937 & 0.0010 & \textbf{Unseen (20)} \\
\bottomrule
\end{tabular}
\end{table}

CRISPRa shows clear improvement for GEARS at 20 epochs. Pearson $\Delta$ top-20 DEGs of
\textbf{0.398} and DA marker delta Pearson of \textbf{0.498}, with sign accuracy of
\textbf{79.3\%} --- the strongest directional prediction across all models and modalities.
The DA marker Pearson delta of 0.498 is particularly notable: the model correctly predicts the
direction and magnitude of perturbation effects on canonical DA marker genes (TH, SLC6A3, NR4A2,
FOXA2) for unseen gene activations with moderate fidelity --- a prerequisite for meaningful DA
identity scoring. Gene activation (CRISPRa) produces larger transcriptional responses than
repression (CRISPRi), making delta metrics more reliable and the task more tractable.

\subsection{Training Dynamics}

GEARS on local hardware (GTX 1650 Ti, 4 GB VRAM) trained for 5 epochs ($\sim$25 minutes) with
validation top-20 DEG MSE improving from 0.0044 to 0.0041, indicating non-convergence. The
20-epoch Colab run confirms convergence: val DE MSE dropped from 0.0062 (epoch 1) to 0.0057
(epoch 9) then plateaued, with best model selected at epoch 9. Final Pearson DE: \textbf{0.9776}.
scGen converged within $\sim$11 epochs (early stopping, patience = 10, ELBO $\sim$900--1200).

\subsection{GEARS 20-Epoch Full-Gene Results (Yumejichi Fujita)}

Yumejichi Fujita ran GEARS independently on GSE152988 CRISPRi with 20 epochs and all 33,538
genes (no HVG subsampling), using a manual train/val/test split of 128/28/28 perturbations.
Due to a library compatibility issue with STRING on the full gene set, an \textbf{empty
co-expression graph} was used.

\begin{table}[h]
\centering
\caption{GEARS 20-epoch (Yumejichi Fujita) --- absolute metrics, CRISPRi.}
\small
\begin{tabular}{ll}
\toprule
Metric & Value \\
\midrule
Pearson $r$ (all genes)    & 0.9986 \\
Pearson $r$ (top-20 DEGs)  & 0.9682 \\
MSE (all genes)            & 0.1000 \\
MSE (top-20 DEGs)          & 0.0488 \\
Train / Val / Test perts   & 128 / 28 / 28 \\
Genes used                 & 33,538 (all) \\
Co-expression graph        & Empty (compatibility bug) \\
Epochs                     & 20 \\
\bottomrule
\end{tabular}
\end{table}

Pearson DE is nearly identical to Katherine's run (0.968 vs.\ 0.973), with the small gap
attributable to fewer training perturbations (128 vs.\ 148) and the missing co-expression graph.
The higher MSE (0.1000 vs.\ 0.00137) is not directly comparable: Yumejichi's MSE is computed
cell-level on all 33k genes while Katherine's is on mean expression vectors. The close Pearson DE
despite the empty co-expression graph suggests STRING provides limited additional signal for
absolute correlation on this dataset.

\subsection{STATE/CPA --- Combined Dataset Results (Zihan Jin, in progress)}

Zihan Jin is fine-tuning STATE on the combined GSE152988 dataset (CRISPRi + CRISPRa merged),
using modality-specific labels (\texttt{CRISPRI::GENE}, \texttt{CRISPRA::GENE}) and a
228/28/28 perturbation split. The base model is \texttt{ST-SE-Replogle}, pretrained on the
Replogle et al.\ 2022 genome-wide perturbation screen. Zihan has additionally trained STATE on
GSE124703 (both iPSC and Day 7 neuron stages), enabling zero-shot cell type evaluation.

\begin{table}[h]
\centering
\small
\begin{tabular}{ll}
\toprule
Configuration & Status \\
\midrule
STATE on GSE152988 combined (228/28/28) & In progress \\
STATE on GSE124703 (iPSC + Day 7 neuron) & In progress \\
CPA via STATE on GSE152988 combined & In progress \\
Zero-shot cross-dataset: GSE152988 $\to$ GSE124703 & Planned after training \\
\bottomrule
\end{tabular}
\end{table}

%% ─────────────────────────────────────────────────────────────────────────────
\section{Research Roadmap}

\subsection{Immediate (before final report)}

\textbf{STATE/CPA results (Zihan Jin):} Integrate benchmark results into final comparison
tables alongside GEARS and scGen.

\textbf{Zero-shot cross-dataset evaluation (Track A):} After STATE fine-tuning on GSE152988, evaluate on GSE124703
(zero-shot stage and cell-type transfer) before using GSE124703 in any training.

\subsection{Extended Experiments}

\textbf{DA identity scoring:} (1) Preprocess GSE140231 to extract clean DA neuron labels;
(2) score GSE152988 cells using the 71-gene DA marker panel; (3) rank perturbations by predicted
DA identity shift under each model. This produces the gene priority list for experimental
validation in PD.

\textbf{scFoundation + GEARS (HPC):} Replace GEARS' STRING node features with scFoundation
context-aware gene embeddings (cells $\times$ 19,264 $\times$ 512). Requires 40+ GB VRAM;
planned for university HPC or Colab A100.

\textbf{Data leakage verification:} Verify whether GSE152988/GSE124703 appear in scFoundation's
CellxGene pretraining corpus or STATE's Replogle pretraining data.

\textbf{CRISPRi vs.\ CRISPRa analysis:} Analyze whether models that excel at predicting
knockdowns also generalize to activation, and vice versa.

%% ─────────────────────────────────────────────────────────────────────────────
\section{Conclusion}

We have established a working multi-model benchmark pipeline for AIVC perturbation prediction on
iPSC-derived neuron data relevant to Parkinson's disease. We introduced a two-tier evaluation
framework that separates absolute expression reconstruction (secondary, inflated by cell identity
background) from delta-based perturbation-specific prediction (primary). Under this framework:
GEARS (20 epochs, Colab A100) achieves Pearson DE of 0.978 (CRISPRi) and 0.961 (CRISPRa) on
completely unseen perturbations, with sign accuracy of 62.3\%/79.3\% and DA marker Pearson
delta of \textbf{0.498} on CRISPRa --- demonstrating meaningful predictive fidelity on the gene
panel directly relevant to DA neuron identity. scGen and mean baseline perform well on the easier
seen-perturbation task but cannot generalize to unseen conditions by design. STATE/CPA results
(Zihan Jin) and the Yumejichi Fujita 20-epoch full-gene run (0.968 Pearson DE) will be integrated
in the final report.

%% ─────────────────────────────────────────────────────────────────────────────
\bibliographystyle{unsrtnat}
\begin{thebibliography}{9}

\bibitem{lopez2018}
Lopez, R. et al.\ (2018).
Deep generative modeling for single-cell transcriptomics.
\textit{Nature Methods}, 15, 1053--1058.

\bibitem{roohani2023}
Roohani, Y. et al.\ (2023).
Predicting transcriptional outcomes of novel multigene perturbations with GEARS.
\textit{Nature Biotechnology}, 42, 927--935.

\bibitem{hao2024}
Hao, M. et al.\ (2024).
Large-scale foundation model on single-cell transcriptomics.
\textit{Nature Methods}, 21, 1481--1491.

\bibitem{lotfollahi2021}
Lotfollahi, M. et al.\ (2021).
Compositional perturbation autoencoder for single-cell response modeling.
\textit{bioRxiv}.

\bibitem{cui2024}
Cui, H. et al.\ (2024).
scGPT: toward building a foundation model for single-cell multi-omics using generative AI.
\textit{Nature Methods}, 21, 1470--1480.

\bibitem{tian2021}
Tian, R. \& Kampmann, M. (2021).
Genome-wide CRISPRi/a screens in human neurons link lysosomal failure to ferroptosis.
\textit{Nature Neuroscience}, 24, 1020--1034.

\bibitem{tian2019}
Tian, R. et al.\ (2019).
CRISPR interference-based platform for multimodal genetic screens in human iPSC-derived neurons.
\textit{Neuron}, 104, 239--255.

\bibitem{boiarsky2023}
Boiarsky, R. et al.\ (2023).
A single cell foundation model.
\textit{bioRxiv}.

\bibitem{wenk2024}
Wenk, P. et al.\ (2024).
Simple baselines for perturbation modelling in single-cell data.
\textit{bioRxiv}.

\end{thebibliography}

\vspace{6pt}
{\small Code and scripts: \url{https://github.com/Katherine-Deborah/AIVC-DA-perturbation}}

\end{document}
